% !TEX root = ../../cookbook.tex
\newpage
\recipe{\textit{Cooked Wine} Tart}{Made with Swiss ``\textit{Vin Cuit}'' (literally ``cooked wine''). A.k.a. ``Tarte au \textit{Raisinée}.''}{}{}

\subsection*{Ingredients}
\begin{ingredients}
	\item 250 g sweet pastry (see recipe on page \pageref{bundner_nusstorte})
	\item 3 eggs
	\item 1 egg yolk
	\item 1.4 dl double cream\footnote{A rich Swiss cream with about 45\% fat.}
	\item 1 dl Vaud-style \textit{cooked wine}
\end{ingredients}

\medskip

% --- Preparation ---
Prepare the \textbf{sweet pastry} and leave it to rest.  
Roll out the pastry to a thickness of \sfrac{1}{2} cm and line a buttered and floured 20 cm diameter tart tin.  
Blind bake the pastry\footnote{To prevent the pastry base from puffing up during blind baking, cover it with a disc of aluminium foil or baking paper and fill with chickpeas to keep it in place.}, for approximately 25 minutes, in an oven at 250°C.  
Allow to cool without removing it from the tin.  

Prepare the custard filling by whisking together the \textbf{remaining ingredients}\footnote{Depending on the degree of reduction of the cooked wine, the quantity may need to be increased up to 1.2 dl.}.  
Fill the pastry shell, place it in the preheated oven at 180°C, and finish filling it to the brim once the tart is in the oven, using a spoon.

Leave the oven door slightly ajar until the custard has set, approximately half an hour, and check the degree of doneness by gently shaking the tin.  
Remove from the tin while still lukewarm and place on a wire rack.

\vspace*{0.75 cm}

% --- Description ---
\begin{recipeblock}
	\noindent 
\end{recipeblock}
